\documentclass[11pt,letterpaper]{article}

\usepackage[T1]{fontenc}
\usepackage{amsmath,amssymb}
\usepackage{tikz}
\usepackage{pgfplots}
\pgfplotsset{compat=1.18}
\usepackage{algorithm}
\usepackage{algpseudocode}
\usepackage{booktabs}
\usepackage{enumitem}
\usepackage[hidelinks]{hyperref}

\newcommand{\N}{\mathbb{N}}
\newcommand{\Z}{\mathbb{Z}}
\newcommand{\Q}{\mathbb{Q}}
\newcommand{\R}{\mathbb{R}}
\newcommand{\C}{\mathbb{C}}
\newcommand{\Map}[3]{#1 \colon #2 \to #3}
\newcommand{\Mod}[1]{#1\text{-}\mathrm{Mod}}
\newcommand{\GroupHom}{\mathrm{Hom}_{\mathbf{Grp}}}
\newcommand{\RingHom}{\mathrm{Hom}_{\mathbf{Ring}}}

\newtheorem{theorem}{Theorem}
\newtheorem{lemma}[theorem]{Lemma}
\newtheorem{corollary}[theorem]{Corollary}
\newtheorem{definition}[theorem]{Definition}
\newtheorem{remark}[theorem]{Remark}
\newenvironment{proof}{\noindent\textit{Proof.}}{\hfill\(\square\)}

\title{Generalized Collatz Library}
\author{}
\date{\today}

\begin{document}
\maketitle

\begin{abstract}
Two affine returns with offset \(1\) cannot close above state \(1\) when
 their total number of binary divisions is even. Substitution places the
multiplier below an integer square root, while a product comparison forces
an incompatible lower bound. The algebraic argument needs no parity
assumption on the multiplier or states. Complementary factors also recover
the multiplier from their sum and product and show that
\(2^{2m}-q^2=7\) has only \((m,q)=(2,3)\) in natural numbers with
\(q\le2^m\). The state-\(1\) cycle explains the sharp lower bound.
\end{abstract}

\section{Two returns and their division total}

For an odd positive multiplier \(q\), consider the raw map
\[
C_q(n)=
\begin{cases}
n/2,&n\text{ even},\\
qn+1,&n\text{ odd}.
\end{cases}
\]
If two successive odd returns close a cycle, their states \(x,y\) and
division counts \(a,b\) satisfy
\begin{equation}
\label{eq:letters}
2^a y=qx+1,\qquad 2^b x=qy+1.
\end{equation}
The states may coincide. The total \(K=a+b\) counts divisions; the two
returns together use \(K+2\) raw ticks. For odd \(q,x,y\), the equations
imply \(a,b\ge1\). The following algebraic lemmas allow arbitrary natural
\(q,x,y,a,b\), so they also cover the original return equations.

\begin{lemma}[Closure and square window]
\label{lem:closure}
The equations \eqref{eq:letters} imply
\[
2^{a+b}x=q^2x+q+2^a
\quad\text{and}\quad
q^2<2^{a+b}.
\]
\end{lemma}

\begin{proof}
Substitution gives
\begin{align*}
2^{a+b}x
&=2^a(2^b x)=2^a(qy+1)\\
&=q(2^a y)+2^a=q(qx+1)+2^a\\
&=q^2x+q+2^a.
\end{align*}
The offset \(q+2^a\) is positive. Thus \(q^2x<2^{a+b}x\).
If \(2^{a+b}\le q^2\), multiplication by \(x\ge0\) contradicts
that strict inequality. This proves the window.
\end{proof}

\begin{lemma}[Product comparison]
\label{lem:product}
Suppose \eqref{eq:letters} holds and \(\mu\le x,y\). Then
\[
\mu^2 2^{a+b}\le(\mu q+1)^2.
\]
\end{lemma}

\begin{proof}
The second equation forces \(x>0\), since \(qy+1>0\); the first likewise
forces \(y>0\). The hypothesis \(\mu\le x\) gives the exact slack identity
\[
\mu(qx+1)+(x-\mu)=(\mu q+1)x.
\]
Indeed, expanding the left side and using \(\mu+(x-\mu)=x\) leaves
\(\mu qx+x\). Dropping the nonnegative slack and using the first equation
therefore gives \(\mu2^a y\le(\mu q+1)x\). Similarly,
\(\mu2^b x\le(\mu q+1)y\). Multiplying yields
\[
\mu^2 2^{a+b}xy\le(\mu q+1)^2xy.
\]
Cancel the positive product \(xy\).
\end{proof}

\begin{lemma}[Strict integer-root ceiling]
\label{lem:ceiling}
If \(q^2<2^{2t}\) and \(\mu\ge2\), then
\[
q\le2^t-1,
\qquad
\mu q+1\le\mu2^t-1,
\qquad
(\mu q+1)^2<\mu^2 2^{2t}.
\]
\end{lemma}

\begin{proof}
The identity \((2^t)^2=2^{2t}\) shows that \(q\ge2^t\) would contradict
the square window. Since the parameters are integers, \(q\le2^t-1\).
Consequently,
\[
\mu q+1\le\mu(2^t-1)+1
=\mu2^t-\mu+1\le\mu2^t-1<\mu2^t.
\]
Here \(\mu\ge2\) supplies the penultimate inequality, and
\(\mu2^t>0\). Squaring the strict comparison gives the result.
\end{proof}

\begin{theorem}[Even division totals]
\label{thm:even-empty}
Let all parameters be natural numbers. If \(a+b=2t\) and \(x,y\ge\mu\ge2\),
the equations \eqref{eq:letters} cannot both hold.
\end{theorem}

\begin{proof}
Lemma~\ref{lem:closure} gives \(q^2<2^{2t}\).
Lemma~\ref{lem:ceiling} gives \((\mu q+1)^2<\mu^2 2^{2t}\).
Lemma~\ref{lem:product} gives the reverse weak inequality.
\end{proof}

Thus a two-return cycle of the raw map with even division total must visit
state \(1\). No enumeration of the states, the multiplier, or the split of
\(2t\) is needed. When \(t\ge1\), the integer \(2^t-1\) is odd and
itself lies in the square window, since its square is less than \((2^t)^2\).
The strict ceiling therefore includes the largest possible odd multiplier.

\section{Complementary factors and recovery}

\begin{lemma}[Factor identities]
\label{lem:split}
If \(q\le2^m\), then the natural numbers \(u=2^m-q\), \(v=2^m+q\) satisfy
\[
uv=2^{2m}-q^2,\qquad u+v=2^{m+1},\qquad v-u=2q.
\]
\end{lemma}

\begin{proof}
The product is the difference of squares. For the sum,
\((2^m-q)+q=2^m\), because \(q\le2^m\), so
\(u+v=2\cdot2^m=2^{m+1}\). Finally \(u+2q=v\), giving the gap.
\end{proof}

\begin{theorem}[Factor recovery]
\label{thm:recovery}
Suppose \(q^2\le2^{2m}\), \(u+v=2^{m+1}\),
\(uv=2^{2m}-q^2\), and \(u\le v\). Then \(v-u=2q\).
\end{theorem}

\begin{proof}
Write \(R=2^m\). The sum is \(u+v=2R\). If \(v<R\), then
\(u\le v<R\) makes the sum less than \(2R\), a contradiction.
Thus \(v\ge R\) and \(u=2R-v\le R\).
Set \(q'=R-u\). Then \(0\le q'\le R\), \(u=R-q'\), and
\(v=R+q'\). Lemma~\ref{lem:split} gives
\(uv=R^2-(q')^2\). Both \(q^2\) and \((q')^2\) are at most \(R^2\),
so restoring these squares to their respective difference identities gives
\(q^2=(q')^2\). The square function is strictly increasing on natural
numbers, hence \(q=q'\). The factor gap in Lemma~\ref{lem:split} is
therefore \(v-u=2q'=2q\).
\end{proof}

The identities and recovery include \(q=0\), when the factors are equal,
and \(q=2^m\), when the smaller factor is zero.

\begin{corollary}[Difference seven]
\label{cor:seven}
If \(q\le2^m\) and \(2^{2m}-q^2=7\), then \(m=2\) and \(q=3\).
Conversely, this pair satisfies the equation.
\end{corollary}

\begin{proof}
Lemma~\ref{lem:split} gives an ordered positive factor pair \(u\le v\)
with \(uv=7\). If \(u\ge3\), then \(uv\ge u^2\ge9\), a contradiction.
If \(u=2\), the product would be even. Thus \(u=1\), \(v=7\).
The factor sum is \(2^{m+1}=8\). The powers of two are strictly increasing,
and \(2^3=8\), so \(m+1=3\) and \(m=2\). The gap gives
\(2q=7-1=6\), hence \(q=3\). Conversely, \(2^4-3^2=16-9=7\).
\end{proof}

\section{Sharpness and exact checks}

At \(q=3\), \(x=y=1\), \(a=b=2\), both return equations read
\(4=3+1\). The product comparison at minimum \(1\) is equality:
\((1\cdot3+1)^2=2^4=16\). The corresponding comparison at minimum \(2\)
is strict: \(7^2=49<4\cdot2^4=64\). This is precisely the state excluded
by the lower bound in Theorem~\ref{thm:even-empty}.

The Lean file \texttt{ReverseBound.lean} proves all algebraic statements,
including the unrestricted multiplier theorem and its raw-map specialisation.
The Go programme checks ranks \(1\le t,m\le12\), every multiplier up to
 the appropriate root, and minima \(2\le\mu\le8\). It also checks factor
recovery, nonnegative slack, and present closure equations in the grid
\(0\le a,b\le4\), \(0\le q\le8\), \(1\le x,y\le8\).
Every operation is checked in unsigned 64-bit arithmetic. Unsupported powers
and overflow are rejected, including exponent wraparound. Independent rank
checks use disjoint worker slots; the numerical loops allocate no heap memory.
The finite checks support the implementation; Theorem~\ref{thm:even-empty}
and the recovery statements have no numerical upper bound.

\end{document}
